\section{Related Work}
\label{sec:related-work}

\textbf{Autonomous Agents.}
The pursuit of creating autonomous agents that can operate intelligently in real-world environments without human involvement has been a persistent goal throughout the history of AI \citep{wooldridge1995intelligent,Minsky88,bubeck2023sparks}. Recently LLMs \citep{touvron2023llama,DBLP:journals/corr/abs-2303-08774} have opened up new opportunities to achieve this goal. These LLMs possess remarkable understanding, reasoning, and generation capabilities, allowing autonomous agents to utilize them as a backbone for handling increasingly complex scenarios~\citep{richards2023auto,nakajima2023babyagi,agentgpt,liu2023agentbench}. However, even though these autonomous agents already demonstrate considerable power, they still lack certain essential human-analogous cognitive capabilities. Hence, some research designs external mechanisms that endow agents with reflection \citep{yao2023react,shinn2023reflexion}, task decomposition \citep{wei2023chainofthought,yao2023tree}, and tool utilization/creation \citep{schick2023toolformer,qin2023tool,qin2023toolllm,qian2023creator} capabilities, which bring autonomous agents closer to achieving artificial general intelligence.



\textbf{Multi-agent System.}
In human society, a well-organized group composed of individual humans can often collaboratively handle a greater workload and accomplish complex tasks with higher efficiency and effectiveness. In the field of AI, researchers draw inspiration from human society and aim to enhance work efficiency and effectiveness by leveraging cooperation among individuals through the study of multi-agent systems (MAS) \citep{multiagent_system}, also referred to as a \textit{multi-agent group} in this paper. The multi-agent group collaboratively makes decisions and executes corresponding actions in a distributed and parallel manner to achieve the common goal, which significantly improves work efficiency and effectiveness. Previous works have leveraged multi-agent joint training to achieve this goal. Recently, some studies have attempted to leverage the intelligence and capabilities of agents for autonomous collaboration. \cite{DBLP:journals/corr/abs-2303-17760} have conceptualized assemblies of agents as a group, and focused on exploring the potential of their cooperation. \cite{DBLP:journals/corr/abs-2304-03442} found social behaviors autonomously emerge within a group of agents, and \cite{DuMultiagentDebate,wang2023unleashing,zhang2023building,DBLP:journals/corr/abs-2307-07924,chan2023chateval} further leverage multi-agent cooperation to achieve better performance on reasoning tasks. Based on these findings, we introduce a framework, denoted as \textit{\framework}, capable of leveraging group cooperation to manage more intricate scenarios. This framework can dynamically adjust its composition according to the current state, aiming to facilitate optimal decision-making and execution.




